% GENERATED FILE. Edit canonical lessons, then run npm run generate:books.
% canonical-source-hash: fnv1a64:6e207f177b1c51c9
% canonical-lessons: BN-C192-hotash, BN-C192-sorol, BN-C192-jotil, BN-C192-alada, BN-C192-eki

\chapter{Disappointed, Simple, Complicated, Separate, Same}
\label{ch:bn-disappointed-simple-complicated-separate-same}
\hlchaptermodality{192}

\begin{chapteropening}
\textbf{By the end of this chapter:} \emph{I can say disappointed, simple, complicated, separate and the same.}

Complete the last lesson of chapter 192: I can say disappointed, simple, complicated, separate and the same.
\end{chapteropening}

\section[\texorpdfstring{\bn{হতাশ} (hôtāsh)}{হতাশ (hôtāsh)}]{\bn{হতাশ} (hôtāsh) — disappointed}

\label{lesson:BN-C192-hotash}



\begin{quote}
Before the new one: say the Bengali for ancient, then the Bengali for real.
\end{quote}

\subsection*{You'll want to know: \bn{হতাশ}}
\textbf{\bn{হতাশ}} — \emph{hôtāsh} — \textquotedblleft{}disappointed\textquotedblright{}.

\begin{grammarlens}[title={What you've built}]
One more word for this chapter.
\end{grammarlens}

\subsection*{Your turn}
\begin{itemize}
\raggedright
  \item \emph{Say it:} \emph{hôtāsh}
  \item \emph{Say it:} \emph{hôtāsh}, once more
  \item \emph{Say it:} say \emph{hôtāsh}
  \item \emph{Recall:} say the Bengali for ancient, then the Bengali for real, then say \emph{hôtāsh} again
\end{itemize}

\begin{tcolorbox}[breakable,colback=teal!4,colframe=teal!35!black,title={Before you move on}]
What does \bn{হতাশ} mean? (\textquotedblleft{}Disappointed\textquotedblright{}.) Say it once more.
\end{tcolorbox}
\section[\texorpdfstring{\bn{সরল} (sôrôl)}{সরল (sôrôl)}]{\bn{সরল} (sôrôl) — simple, innocent}

\label{lesson:BN-C192-sorol}



\begin{quote}
Before the new one: say the Bengali for real, then the Bengali for disappointed.
\end{quote}

\subsection*{You'll want to know: \bn{সরল}}
\textbf{\bn{সরল}} — \emph{sôrôl} — \textquotedblleft{}simple, innocent\textquotedblright{}.

\begin{grammarlens}[title={What you've built}]
One more word for this chapter.
\end{grammarlens}

\subsection*{Your turn}
\begin{itemize}
\raggedright
  \item \emph{Say it:} \emph{sôrôl}
  \item \emph{Say it:} \emph{sôrôl}, once more
  \item \emph{Say it:} say \emph{sôrôl}
  \item \emph{Recall:} say the Bengali for real, then the Bengali for disappointed, then say \emph{sôrôl} again
\end{itemize}

\begin{tcolorbox}[breakable,colback=teal!4,colframe=teal!35!black,title={Before you move on}]
What does \bn{সরল} mean? (\textquotedblleft{}Simple, innocent\textquotedblright{}.) Say it once more.
\end{tcolorbox}
\section[\texorpdfstring{\bn{জটিল} (jôṭil)}{জটিল (jôṭil)}]{\bn{জটিল} (jôṭil) — complicated}

\label{lesson:BN-C192-jotil}



\begin{quote}
Before the new one: say the Bengali for disappointed, then the Bengali for simple.
\end{quote}

\subsection*{You'll want to know: \bn{জটিল}}
\textbf{\bn{জটিল}} — \emph{jôṭil} — \textquotedblleft{}complicated\textquotedblright{}.

\begin{grammarlens}[title={What you've built}]
One more word for this chapter.
\end{grammarlens}

\subsection*{Your turn}
\begin{itemize}
\raggedright
  \item \emph{Say it:} \emph{jôṭil}
  \item \emph{Say it:} \emph{jôṭil}, once more
  \item \emph{Say it:} say \emph{jôṭil}
  \item \emph{Recall:} say the Bengali for disappointed, then the Bengali for simple, then say \emph{jôṭil} again
\end{itemize}

\begin{tcolorbox}[breakable,colback=teal!4,colframe=teal!35!black,title={Before you move on}]
What does \bn{জটিল} mean? (\textquotedblleft{}Complicated\textquotedblright{}.) Say it once more.
\end{tcolorbox}
\section[\texorpdfstring{\bn{আলাদা} (ālādā)}{আলাদা (ālādā)}]{\bn{আলাদা} (ālādā) — separate, different}

\label{lesson:BN-C192-alada}



\begin{quote}
Before the new one: say the Bengali for simple, then the Bengali for complicated.
\end{quote}

\subsection*{You'll want to know: \bn{আলাদা}}
\textbf{\bn{আলাদা}} — \emph{ālādā} — \textquotedblleft{}separate, different\textquotedblright{}.

\begin{grammarlens}[title={What you've built}]
One more word for this chapter.
\end{grammarlens}

\subsection*{Your turn}
\begin{itemize}
\raggedright
  \item \emph{Say it:} \emph{ālādā}
  \item \emph{Say it:} \emph{ālādā}, once more
  \item \emph{Say it:} say \emph{ālādā}
  \item \emph{Recall:} say the Bengali for simple, then the Bengali for complicated, then say \emph{ālādā} again
\end{itemize}

\begin{tcolorbox}[breakable,colback=teal!4,colframe=teal!35!black,title={Before you move on}]
What does \bn{আলাদা} mean? (\textquotedblleft{}Separate, different\textquotedblright{}.) Say it once more.
\end{tcolorbox}
\section[\texorpdfstring{\bn{একই} (eki)}{একই (eki)}]{\bn{একই} (eki) — the same}

\label{lesson:BN-C192-eki}



\begin{quote}
Before the new one: say the Bengali for complicated, then the Bengali for separate.
\end{quote}

\subsection*{You'll want to know: \bn{একই}}
\textbf{\bn{একই}} — \emph{eki} — \textquotedblleft{}the same\textquotedblright{}.

\begin{grammarlens}[title={What you've built}]
One more word for this chapter.
\end{grammarlens}

\subsection*{Your turn}
\begin{itemize}
\raggedright
  \item \emph{Say it:} \emph{eki}
  \item \emph{Say it:} \emph{eki}, once more
  \item \emph{Say it:} say \emph{eki}
  \item \emph{Recall:} say the Bengali for complicated, then the Bengali for separate, then say \emph{eki} again
\end{itemize}

\begin{tcolorbox}[breakable,colback=teal!4,colframe=teal!35!black,title={Before you move on}]
What does \bn{একই} mean? (\textquotedblleft{}The same\textquotedblright{}.) Say it once more.
\end{tcolorbox}
